
%%%%%%%%%%%%%%%%%%%%%%%%%%%%%%%%%%%%%%%%%%%%%%%%%%%%%%%%%%%%%
\subsubsection{具体分析}

问题二要求在超低排放约束（$C_{out}\leq 10\;\text{mg/Nm}^3$）下，寻求各电场操作参数的最优组合以最小化总电耗，属于带约束的非线性优化类问题。

\textbf{工况划分的必要性。}入口条件$C_{in}\in[18,72]\;\text{g/Nm}^3$、$T_{in}\in[111.7,158.2]\;^\circ\text{C}$、$Q$波动范围大，用一个参数组合覆盖所有工况不现实——高浓度工况需要更高电压和更短振打周期，低浓度工况则可降低电压省电。因此需将连续工况空间离散化为若干典型场景，在每个场景内独立寻优。

\textbf{工况划分方法选择。}等频分箱在多维空间易产生空格或不均衡区间，而K-Means在$(C_{in},T_{in},Q)$三维标准化空间内自动寻找方差最小的聚团，符合"区间内相近、区间间差异明显"的工程直觉。聚类数$K$通过轮廓系数在$K\in[2,8]$范围内选取，同时辅以肘部法则观察组内平方和随$K$的下降拐点，双验证确定$K=6$。校验每工况样本数$\geq 3\%$总样本以保证统计可靠性。三个特征量纲差异巨大（$C_{in}$约36、$T_{in}$约126、$Q$约462830），需先\texttt{StandardScaler}去量纲再聚类，否则$Q$会主导聚类结果。

\textbf{温度作为聚类特征但不进入效率模型的说明。}$T_{in}$参与工况划分但不直接出现在Deutsch效率模型中。温度对除尘效率的影响主要通过粉尘比电阻（高温低比电阻区、低温高比电阻区）间接作用于驱进速度$\omega$，而比电阻效应已隐含在$kA_0$的拟合值中（$kA_0$是全数据集的统计平均，已吸收了温度对比电阻的平均影响）。将$T_{in}$作为聚类特征的理由是：不同温度区间的入口条件组合不同（高温往往伴随低浓度——余热锅炉投运时烟气降温增湿），温度帮助K-Means更准确地区分工况边界，即使效率模型中不显式含温度项，工况划分的准确性仍受益于温度维度。

\textbf{温度特征的聚类增益验证。}为量化$T_{in}$对聚类的贡献，对比$(C_{in},Q)$二维聚类与$(C_{in},T_{in},Q)$三维聚类的轮廓系数：在$K=6$时，三维轮廓系数0.353略低于二维0.394，但三维聚类将同浓度不同温度的样本（如$C_{in}\approx44$但$T_{in}=119$℃与$131$℃）划分为不同工况，使各工况内温度方差更小、入口条件更均匀，有利于后续寻优时用工况均值代表整个区间。若去掉$T_{in}$，同浓度不同温度的样本被合并，工况内温度跨度增大，优化解对温度极端子样本的代表性下降。

\textbf{约束条件梳理。}优化需满足以下约束：（1）排放约束$C_{out}\leq C_{limit}=10\;\text{mg/Nm}^3$；（2）电压上下限$U_{min,i}\leq U_i\leq U_{max,i}$，其中$U_{min,i}$取历史最小值，$U_{max,i}$取历史最大值的1.2倍（变压器允许短时过载20\%，ESP高压整流变压器通常按额定容量120\%设计短时过载能力）；（3）振打周期上下限$T_{min,i}\leq T_i\leq\min(T_{max,i},T_{crit,i})$，其中$T_{crit,i}$取历史95分位数作为安全上限，防止积灰退化。

各电场优化边界具体数值如表\ref{tab:优化边界}所示。

\begin{table}[H]
  \centering
  \caption{各电场操作参数优化边界}
  \label{tab:优化边界}
  \renewcommand{\arraystretch}{1.1}
  \begin{tabularx}{\textwidth}{CCCCCCC}
    \toprule
    电场 & $U_{min}$(kV) & $U_{max}$(kV) & 历史最大值(kV) & $T_{min}$(s) & $T_{max}$(s) & $T_{crit}$(s) \\
    \midrule
    1 & 45.9 & 85.1 & 70.9 & 157 & 290 & 261 \\
    2 & 45.4 & 85.7 & 71.4 & 154 & 289 & 261 \\
    3 & 38.0 & 71.0 & 59.2 & 359 & 506 & 472 \\
    4 & 37.0 & 69.8 & 58.2 & 358 & 510 & 472 \\
    \bottomrule
  \end{tabularx}
\end{table}

说明：$U_{max}$超出历史最大值约20\%，是变压器短时过载裕度。历史运行中电压未达到此上限，可能是因为（1）历史操作偏保守，未充分利用变压器容量；（2）$C_{out}$传感器限幅，操作人员无法观测到电压提升对排放的改善效果，缺乏提压激励。优化结果中$U_1$贴上界85.1\,kV，说明在超低排放约束下，前电场需要充分提升电压以降低浓度，这与"历史运行未充分利用变压器容量"的判断一致。
